\paragraph{}
A utilização de aprendizagem automática na predição médica cria dois desafios relacionados: os fluxos preditivos publicados podem ser difíceis de reproduzir e as explicações \textit{post-hoc} podem variar consoante o modelo, a representação dos dados e o método de explicação. Esta dissertação desenvolveu e avaliou uma estrutura reprodutível de Inteligência Artificial Explicável (XAI), através da replicação e extensão de cinco estudos com dados médicos relativos a doença cardíaca, acidente vascular cerebral (AVC), diabetes, doença renal crónica e readmissão hospitalar associada à medição de HbA1c.

Os fluxos dos estudos de origem foram reconstruídos em notebooks específicos para cada conjunto de dados. A replicação segundo os procedimentos descritos nos artigos foi distinguida de ramos analíticos corrigidos sempre que a ordem do pré-processamento, a fuga de informação, o desequilíbrio entre classes ou o desenho da avaliação limitavam a validade de uma imitação direta. Quatro modelos finais foram avaliados em dados de teste reservados; o estudo de HbA1c manteve o seu desenho inferencial original aplicado à coorte completa. Uma estrutura comum gerou importância nativa do modelo, quando disponível, importância por permutação, SHAP, explicações LIME repetidas para até dez observações por categoria e 20 sementes de perturbação, medidas de concordância entre ordenações e resultados de estabilidade após reajustamentos em amostras bootstrap quando o retreino estava disponível.

A concordância entre os métodos globais dependeu fortemente do conjunto de dados e da respetiva representação. A importância nativa e o SHAP apresentaram elevada concordância nos modelos de doença cardíaca, AVC, diabetes e doença renal crónica, enquanto as ordenações obtidas por permutação divergiram com maior frequência. No modelo de HbA1c, rico em interações, o SHAP e a importância por permutação apresentaram elevada concordância, enquanto a magnitude absoluta dos coeficientes mostrou uma concordância quase nula com o SHAP. Ao longo de 100 reajustamentos bootstrap, a estabilidade da ordenação completa foi mais elevada para doença renal crónica, seguida de diabetes e AVC, enquanto a doença cardíaca apresentou uma persistência inferior, embora positiva. A diabetes apresentou a maior sobreposição entre as três variáveis principais. As correlações médias entre ordenações LIME obtidas com diferentes sementes variaram entre 0,898 e 0,969, mas previsões corretas e incorretas partilhavam frequentemente as mesmas regras locais dominantes, incluindo erros de AVC com elevada confiança. Consequentemente, explicações coerentes e repetíveis não estabeleceram, por si só, a correção das previsões.

A estrutura foi executada com sucesso em classificadores lineares, modelos de conjunto baseados em árvores e na implementação de regressão logística testada. Foi apoiada por versões fixas das dependências, metadados de proveniência, pacotes de modelos guardados, reexecuções através da linha de comandos, relatórios de comparação e testes automatizados. A sua generalisabilidade permanece limitada às interfaces de classificação binária avaliadas. Não foram realizadas avaliação prospetiva, análise causal, avaliação centrada nos profissionais de saúde ou validação da utilidade clínica. A principal conclusão é que a XAI deve ser tratada como um processo reprodutível e multimétodo de auditoria de modelos, e não como um certificado automático de confiança clínica.

\paragraph{}
\textbf{Palavras-chave:} Inteligência Artificial Explicável; Predição Clínica; Reprodutibilidade; SHAP; LIME; Importância de Variáveis; Estabilidade das Explicações.
